\section{Background and Threat Model}\label{sec:threat}

\paragraph{NLIP} An NLIP message (ECMA-430) is a JSON object with a first submessage
(\emph{format}, \emph{subformat}, \emph{content}) and an optional ordered list of further
submessages. Reserved \emph{token} submessages carry control data such as the conversation
token and an authorization credential. ECMA-431 binds NLIP to HTTP; the course framework serves it at \texttt{POST /nlip}.
The course framework runs pluggable SEs at the server's ingress and egress, each returning
allow, block (HTTP~403) or replace. The course's architecture defines a Type-3 SE as
one that secures the agent's interactions with tools and models and may be invoked directly
before or after a tool call; the framework itself has hooks only at ingress and egress.

\paragraph{Attacker.} The attacker controls (i)~the \emph{metadata} of a tool the agent
loads from an MCP server: the tool description and the descriptions of its parameters
(\texttt{inputSchema}), which the model reads but the user does not~\cite{invariant2025,mcptox};
and (ii)~\emph{content} that tools return: an email body, a calendar event or a shared
file~\cite{greshake2023,agentdojo}. The attacker controls only this metadata and content, for example by pushing a malicious
description update through a registry or proxy; in our tests the poison sits in the
descriptions of the workspace server's own tools. An attacker who controls the server's
backend could act without the agent and is out of scope. Poison in one server that steers calls
to another (shadowing) was not tested; a newly added server's tools are blocked as unknown
until the deployer labels them, unless they reuse a labelled name. The attacker does not
control the user, the NLIP client,
the agent host, the models or the guard, and cannot forge the user's Google ID token. The
attacker's goal is a side effect the user did not ask for. AgentDojo's 14 workspace goals
include sending a security code or the inbox to an outside address, deleting files,
inviting an outsider to a meeting, and mailing phishing links.

\paragraph{Why it matters for agents.} A web server that receives an injected string stores
it. An agent reads it, reasons about it, and can act on it with every permission the user
granted, at machine speed and out of the user's sight. A successful injection therefore
becomes an action: mail sent from the user's account, a file deleted, a code exfiltrated.
This is a confused deputy (ECMA-434~\S7.2.3). One poisoned MCP server can affect every user
of every agent that loads it.

\paragraph{Security objectives.} Table~\ref{tab:sco} maps \sys to the Security
Countermeasures and Objectives (SCOs) of ECMA-434~\cite{ecma434}, Clause~7. The primary SCO
is \S7.3.1, \emph{Prevention of Prompt Injection (including Indirect Prompt Injection)},
which names tools as a source of untrusted content. Our end-to-end evidence, however, is for
blocking the actions an injection induces (\S7.3.2, \S7.2.3): the descriptions are not
filtered, and the result scanner, the filtering part, was evaluated only offline. We address individual SCOs and do not
claim conformance to a security profile, which also requires all of Profile~1, including TLS
(\S7.1.1).

\begin{table}[t]
\caption{ECMA-434 SCOs addressed (\S7.x.y, title as in the standard).}\label{tab:sco}
\small
\begin{tabular}{@{}>{\raggedright\arraybackslash}p{0.44\columnwidth}>{\raggedright\arraybackslash}p{0.52\columnwidth}@{}}
\toprule
SCO & How \\
\midrule
\S7.2.1 Identity \& Authentication: NLIP Client Authentication
  & The app's Google OAuth ingress SE (not the guard). \\
\S7.2.2 Runtime \& Behaviour: NLIP Client Authorization
  & Partly: the sign-in SE's optional allow-list admits only listed users; the guard adds task
    alignment, not per-client permissions. \\
\S7.2.3 Runtime \& Behaviour: Privilege Misuse Prevention
  & Calls the user would not approve are refused, even though the agent holds the privilege. \\
\S7.2.4 Observability \& Audit: Agent Behaviours
  & Partly: every decision, with its score and reason, goes to an audit log (single agent; no
    cross-agent tracing or budgets). \\
\S7.3.1 Prevention of Prompt Injection (including Indirect Prompt Injection)
  & Primary threat; partly met. Tool results are scanned and injected spans redacted (no
    end-to-end evaluation); descriptions are not filtered. \\
\S7.3.2 Runtime Agent Behaviour Controls
  & Each side-effecting call is validated against the user's request before it runs; this
    is what the end-to-end results measure. \\
\bottomrule
\end{tabular}
\end{table}
